\section*{Outline and motivation}

In the previous lecture we discussed the use of ray theory in determining the Earth's spherical velocity structure. 
We now consider an extension of these ideas, known as \textbf{delay time tomography}, that can be used to produce quantitative 
images of lateral velocity variations within the Earth. This method falls into a class of techniques known 
as \textbf{seismic tomography}, and we will say more about some other variants in later lectures. Such work is 
done to better understand convection in the mantle, with the tomographic models obtained acting as an indirect 
observational constraint on this process. Delay time tomography is an example of a \textbf{linear inverse problem}, and 
we discuss a practical method for its solution known as \textbf{regularised least squares}.

\section*{Laterally varying structure from seismology}

Though the Earth is close to being spherically symmetric, the small deviations from sphericity are still of great 
interest. This is because they can tell us about processes such as mantle convection deep within the Earth that are 
inaccessible to direct observation. The area of seismology concerned with imaging lateral variations in the Earth's 
interior is known as seismic tomography, and is one of the main areas of current research within solid Earth geophysics. 
The results of such studies are tomographic models that show variations in seismic wave speed and other parameters within 
the Earth. These variations are defined \textbf{relative to a fixed spherically symmetric reference model}, and so include only 
the \textbf{laterally varying perturbations}. An early example of such a tomographic model is shown in Fig.~\ref{fig:1}.

\begin{center}
    \includegraphics[width=0.42\textwidth]{figures/L7F1.png}    
    \captionof{figure}{An image of a tomographic model of S-wave velocity variations within the Earth produced by Woodhouse \& Dziewonski (1989). The colours show the perturbation in S-wave velocity relative to a spherical reference model, with blue areas being
     faster than normal, and red areas slower. Areas of faster velocity are thought to be colder than normal, while those that are
      slower are hotter, and hence the chosen colour scale. These temperature anomalies have associated changes in density, and
       so the models tell us something about the buoyancy forces that drive mantle convection.}
    \label{fig:1}
\end{center}

The basic idea underlying seismic tomography is that if we knew the Earth's internal structure then we could calculate what
would be observed after an earthquake, and so we can try to work backwards to determine this structure from seismic observations.
This is a further example of an inverse problem, and in this lecture we shall start to discuss the solution of such problems in 
detail. In particular, we will focus on travel time measurements that are modelled using ray theory. This was the earliest 
tomographic method developed, with the first papers written in the late 1970s. Many of the basic ideas, however, carry over to 
more complicated tomographic techniques including waveform tomography that we will discuss next time.

\section*{Delay time measurements}

If we know the location and origin time of an earthquake, then we have seen that it is possible to measure the travel times of 
various seismic phases recorded on seismograms around the globe. If the Earth were spherically symmetric then the observed travel 
times would (within the limits of observational errors) depend only on the epicentral angle and the source depth. Though this is 
a reasonable first approximation, there do exist variations in observed travel times due to lateral variations within the Earth. 
For a given travel time measurement, we define its associated \textbf{delay time} by
\begin{equation}
    \delta T = T - T_{0},
    \label{eq:1}
\end{equation}
where $T$ is the \textbf{measured} travel time, and $T_{0}$ is the travel time of the seismic phase \textbf{predicted}
by a spherically symmetric reference model. We note that this definition is such that the delay time will be positive if the observed 
phase arrives later than expected. Such delay times can be measured for many source-receiver pairs and for many seismic 
phases, and constitute the data in delay time tomography. To some extent these times can be measured by hand. But more 
often this is done using signal processing tools such as cross-correlation, or increasingly using forms of machine 
learning.

\section*{Relating delay times to lateral variations}

Suppose we have obtained the delay time of a seismic phase for a given source-receiver pair. Neglecting for the moment 
observational errors, this delay time would vanish if the Earth's velocity structure were equal to that of the spherical 
reference model. In practice, measured delay times are small (e.g.\ tens of seconds for a phase whose travel time is tens 
of minutes), and so we can conclude that the Earth's true structure cannot be too far from the spherical reference model. 
We know that the travel time of a seismic phase can be calculated using ray theory. But this travel time is related to 
the velocity structure in a complicated non-linear manner through solution of the Hamiltonian ray equations, and we 
cannot in general write an explicit formula for $T$ in terms of the velocity, $v$. Because, however, lateral velocity 
variations within the Earth seem to be small we can simplify the problem by seeking an approximate \textbf{linearised relation}
between the observed delay time and velocity perturbations within the Earth. 

Let us write $T(v)$ for the travel time of a certain seismic phase as determined by the velocity model, $v$. This 
is a mapping from the velocity, which is a scalar function in the earth model, to the real numbers, and so the 
mapping $v\mapsto T(v)$ defines a
\textbf{functional}. To proceed, we need to introduce the 
\textbf{functional derivative} of  $T$ with respect to $v$ as defined through a first-order Taylor expansion
\begin{equation}
    T(v+\delta v) = T(v) + \langle DT(v)|\delta v \rangle + O(\delta v^{2}).
    \label{eq:2}
\end{equation}
As in our discussion of functional derivatives in Lecture~13 in the context of Hamilton's principle, we see that this 
functional derivative, $DT(v)$, is an object that acts on a velocity perturbation, $\delta v$, to produce the first-order 
change in the travel time. We denote the action of the functional derivative on the velocity perturbation using the \textbf{inner 
product} notation $\langle \cdot | \cdot \rangle$. 

The desired functional derivative can be obtained very easily using \textbf{Fermat's principle} for elastic ray theory, which you 
established in the first problem set. Having solved the relevant ray tracing equations, we know that the travel time can 
be expressed as the following line integral
\begin{equation}
    T = \int_{\mathrm{ray}} \frac{1}{v} \dd s,
    \label{eq:3}
\end{equation}
where $s$ denotes the arc length along the ray. Fermat's principle states that this travel time is stationary with respect 
to all first-order perturbations to the ray geometry that leave the end-points fixed. If we perturb the velocity of the 
earth model to be $v+\delta v$, then the travel time is affected by both (i) a change in the ray geometry, and (ii) the 
change in velocity along the new ray path. To first-order accuracy the contributions of these effects can be added, but 
by Fermat's principle the first-order term associated with the change in ray geometry is zero. It follows that the delay 
time can be written
\begin{equation}
    \delta T = -\int_{\mathrm{ray}} \frac{\delta v}{v^{2}} \dd s,
    \label{eq:4}
\end{equation}
where this integral is taken along the unperturbed ray. To identify the functional derivative here, it seems natural to use 
the conventional inner product
\begin{equation}
    \langle u|w \rangle = \int_{M} uw \dd^{3}\mathbf{x},
    \label{eq:5}
\end{equation}
for any two real-valued scalar fields defined in the earth model. Looking at eq.~(\ref{eq:4}), it follows that the functional derivative 
can be suitably defined only if we allow it to be a generalised function that is non-zero only along the reference ray path, 
while along the ray path it is infinite in such a way that
\begin{equation}
    \langle DT(v)|\delta v \rangle = -\int_{\mathrm{ray}} \frac{\delta v}{v^{2}} \dd s.
    \label{eq:6}
\end{equation}
This definition should be compared with the familiar one for a Dirac delta function. The singular nature of this functional 
derivative causes difficulties in some approaches to solving the inverse problem, but these issues are avoided (but at a cost) 
within the most common approach discussed below.

\section*{Setting up the inverse problem}

We now consider how the inverse problem can be formulated. Our data in such a study comprise delay times for a variety of seismic 
phases measured from many source-receiver pairs. Typically the number of delay times used will be of order $10^{5}$--$10^{8}$. For 
definiteness, we focus on P-wave phases, meaning the data are only sensitive to variations in the P-wave velocity, $\alpha$. Nearly 
identical methods can be applied to studies based on S-waves, or those that consider both P- and S-wave measurements simultaneously.

At the outset, we shall assume that (i) velocity variations in the Earth are sufficiently small that they can be related linearly 
to the delay times, and (ii) the length scale of heterogeneity in the Earth is such that ray theoretic approximations are 
appropriate. Given these assumptions, each delay time measurement $\delta T_{i}$ for $i=1,\dots,n$ can then be expressed in terms 
of the P-wave velocity perturbation $\delta\alpha$ through the equation
\begin{equation}
    \delta T_{i} = -\int_{\mathrm{ray}_{i}} \frac{\delta\alpha}{\alpha^{2}} \dd s + e_{i},
    \label{eq:7}
\end{equation}
where the integration is taken along the appropriate unperturbed ray path, $\alpha$ denotes the P-wave velocity in the reference model, 
and we have written $e_{i}$ for the associated measurement error. We will assume initially that the $e_{i}$ represent zero-mean, 
uncorrelated Gaussian errors with known variances $\sigma_{i}^{2}$. More complicated statistical models can, of course, be applied, 
but this is not the place for such a discussion.

From the data set $\{\delta T_{i}\}_{i=1}^{n}$ we wish to recover the velocity perturbation $\delta\alpha$, which is a (piecewise) 
continuous function within the earth model. Clearly a finite amount of data can never be sufficient to recover such a continuous 
function uniquely, and so we are dealing with a fundamentally \textbf{underdetermined problem}. There is a range of methods for dealing with 
this problem, but we consider one that is both widely used and comparatively simple (the latter point, of course, contributing to the 
former). First we assume that the velocity perturbation can be expanded using a \textbf{finite number} of basis functions
\begin{equation}
    \delta\alpha(\mathbf{x}) = \sum_{j=1}^{m} \delta\alpha_{j} \varphi_{j}(\mathbf{x}).
    \label{eq:8}
\end{equation}
For example, we might expand the angular dependence of the velocity perturbation in spherical harmonics up to some finite degree, 
and then use a set of polynomial functions to describe the radial dependence. Equally, some studies split the earth model into a 
series of small blocks, assuming a constant velocity in each one. Having expanded the velocity perturbations, the problem is reduced 
to one of linear algebra that can be solved using elementary methods. But clearly we are introducing bias into our solutions because 
we cannot say a priori what an appropriate basis should be.

Accepting this strong assumption on the form of the velocity model, we can then write the $i$th delay time as
\begin{equation}
    \delta T_{i} = \sum_{j=1}^{m} \left(-\int_{\mathrm{ray}_{i}} \frac{\varphi_{j}}{\alpha^{2}} \dd s \right) \delta\alpha_{j} + e_{i},
    \label{eq:9}
\end{equation}
which gives a system of linear equations relating the velocity expansion coefficients $\{\delta\alpha_{j}\}_{j=1}^{m}$ to the delay time data $\{\delta T_{i}\}_{i=1}^{n}$. It will be useful to write this equation as
\begin{equation}
    \mathbf{d} = \mathbf{A}\mathbf{m} + \mathbf{e},
    \label{eq:10}
\end{equation}
where $\mathbf{d}$ is an $n$-dimensional vector of the measured delay times, $\mathbf{m}$ is an $m$-dimensional vector of the 
velocity expansion coefficients, and $\mathbf{e}$ is a vector of the random errors. The elements of the $n$-by-$m$ matrix 
$\mathbf{A}$ are given by
\begin{equation}
    A_{ij} = -\int_{\mathrm{ray}_{i}} \frac{\varphi_{j}}{\alpha^{2}} \dd s,
    \label{eq:11}
\end{equation}
and can be determined from knowledge of the reference model and the basis functions. For the block parameterisation, in 
which $\varphi_{j}$ is equal to one within the $j$th block and to zero elsewhere, $A_{ij}$ is simply minus the length of the 
$i$th ray path within the $j$th block divided by $\alpha^{2}$ (taking the reference velocity to be uniform within a block). 
Each ray crosses only a small fraction of the blocks, and so most of the entries of $\mathbf{A}$ are zero: the matrix is 
\textbf{sparse}. Problems with $10^{8}$ data and $10^{6}$ unknowns are solved using iterative methods that require only 
products of $\mathbf{A}$ and its transpose with vectors, and sparsity makes such products cheap, though it is not essential: 
what matters most is a good iterative method, and that is a topic beyond this course.

To illustrate the ideas of this lecture we will make use of a toy problem, shown in Fig.~\ref{fig:2}, in which straight 
rays cross a square region that has been divided into $24\times24$ blocks (a ``pixel'' basis), the data being the delay 
times along each ray with added Gaussian errors. The model to be recovered is the fractional slowness perturbation in each 
block, and the input model used to generate the data is a smooth random field.

\begin{figure}[ht]
    \centering    
    \includegraphics[width=\textwidth]{figures/L7F2.pdf}        
    \caption{A toy two-dimensional tomographic problem with a pixel basis of $24\times24$ blocks. Stars mark ``earthquakes'', 
    triangles ``seismometers'', and the straight lines denote the ray paths. The colours show the fractional slowness perturbation. 
    On the left is the input model, and on the right the model recovered by regularised least squares with norm damping, the damping 
    parameter having been chosen so that the data are fit to within their errors. In blocks that no ray crosses the recovered model 
    is zero, while within the covered region the anomalies are smeared out along the ray paths.}
\label{fig:2}
\end{figure}

\section*{Non-uniqueness and the null space}
Having reduced delay time tomography to a set of linear algebraic equations
\begin{equation}
    \mathbf{d} = \mathbf{A}\mathbf{m} + \mathbf{e},
    \label{eq:12}
\end{equation}
we can start to think how the model $\mathbf{m}$ might be recovered. If we assume that the observational errors are negligible, 
then we might be tempted to solve this equation using the inverse of $\mathbf{A}$. However, we recall that $\mathbf{A}$ is an 
$n$-by-$m$ matrix, where $n$ is the number of delay time measurements, and $m$ is the number of model parameters, and as we 
will typically have $n \neq m$ the inverse of this matrix cannot be defined. 

We could, of course, select the number of model 
parameters to be equal to the number of data ($n=m$), but there is no physical reason why this should be the case. Indeed, it 
is better to determine the model parameterisation based on some idea of what the data should reasonably be able to resolve 
within the Earth. For example, if we are looking at waves of a given dominant frequency, then we can estimate their wavelength 
in different parts of the Earth using the reference velocity structure. We would not expect to resolve features within the 
Earth whose size is substantially less than the wavelength of the seismic waves, and can build our model parameterisation 
accordingly. When this is done it is often the case that $m \gg n$, which is to say that there are many more model parameters than data.

Considering first the idealised case in which there are no data errors, suppose we have found a solution, $\mathbf{m}$, of the equation
\begin{equation}
    \mathbf{d} = \mathbf{A}\mathbf{m}.
    \label{eq:13}
\end{equation}
If there exists a model vector $\mathbf{m}_{0}$ such that
\begin{equation}
    \mathbf{A}\mathbf{m}_{0} = \mathbf{0},
    \label{eq:14}
\end{equation}
then it is clear that for any scalar $a$ we have
\begin{equation}
    \mathbf{A}(\mathbf{m} + a\mathbf{m}_{0}) = \mathbf{A}\mathbf{m} + a\mathbf{A}\mathbf{m}_{0} = \mathbf{A}\mathbf{m} = \mathbf{d},
    \label{eq:15}
\end{equation}
and so $\mathbf{m} + a\mathbf{m}_{0}$ is still an exact solution to the problem. The set of these vectors forms a linear subspace 
known as the \textbf{null space} of $\mathbf{A}$. If the null space is non-trivial (i.e.\ contains non-zero vectors) then it is clear that the inverse problem cannot 
have a unique solution.

When there are more model parameters than data then $\mathbf{A}$ will always have a non-trivial null space, 
but even if $n \ge m$ there may still be a non-trivial null space due to some of the observations failing to be linearly independent. 
Almost all inverse problems associated with delay time tomography will have a high-dimensional null space, this largely being due 
to insufficient ray coverage throughout the Earth. This is illustrated in Fig.~\ref{fig:3} for the toy problem of Fig.~\ref{fig:2}. 
A model that is non-zero only within cells crossed by no rays obviously lies in the null space, but so do more complicated models 
within the well-covered region, and adding any such model to the input model produces exactly the same data.

\begin{figure}[ht]
    \centering    
    \includegraphics[width=\textwidth]{figures/L7F3.pdf}        
    \caption{A model lying in the null space of the toy tomographic problem shown in Fig.~\ref{fig:2} (left), and the input model of 
    that figure with this null-space model added (right). The two models on the right of Fig.~\ref{fig:2} and here produce 
    identical delay times along every ray, and so cannot be distinguished by the data.}
    \label{fig:3}
\end{figure}


The null space can be determined explicitly through an eigenvalue problem. The $m\times m$ matrix $\mathbf{A}^{T}\mathbf{A}$ 
is symmetric, and hence possesses real eigenvalues and an orthonormal set of eigenvectors $\mathbf{v}_{1},\dots,\mathbf{v}_{m}$. 
Moreover, if $\mathbf{A}^{T}\mathbf{A}\mathbf{v}=\mu\mathbf{v}$ with $\mathbf{v}$ normalised, then 
$\mu=\mathbf{v}^{T}\mathbf{A}^{T}\mathbf{A}\mathbf{v}=\|\mathbf{A}\mathbf{v}\|^{2}\ge0$, and so the eigenvalues can be written 
as $s_{j}^{2}$ with $s_{j}\ge0$,
\begin{equation}
    \mathbf{A}^{T}\mathbf{A}\mathbf{v}_{j}=s_{j}^{2}\mathbf{v}_{j}.
    \label{eq:16}
\end{equation}
The same identity shows that $\mathbf{A}\mathbf{v}_{j}=\mathbf{0}$ if and only if $s_{j}=0$, and hence the null space of 
$\mathbf{A}$ is precisely the span of those eigenvectors having zero eigenvalue. If we expand a model in the eigenvectors 
as $\mathbf{m}=\sum_{j}a_{j}\mathbf{v}_{j}$, the data $\mathbf{d}=\mathbf{A}\mathbf{m}$ depend only on those coefficients with 
$s_{j}>0$, while the remaining coefficients are entirely undetermined. For $s_{j}>0$ we can define the unit vectors 
$\mathbf{u}_{j}=\mathbf{A}\mathbf{v}_{j}/s_{j}$ within the data space, which are orthonormal because 
$\mathbf{u}_{i}\cdot\mathbf{u}_{j}=\mathbf{v}_{i}^{T}\mathbf{A}^{T}\mathbf{A}\mathbf{v}_{j}/(s_{i}s_{j})=\delta_{ij}$, and then write
\begin{equation}
    \mathbf{A}=\sum_{s_{j}>0}s_{j}\mathbf{u}_{j}\mathbf{v}_{j}^{T},
    \label{eq:17}
\end{equation}
which is known as the \textbf{singular value decomposition} of $\mathbf{A}$, with the $s_{j}$ its \textbf{singular values}. 
This decomposition tells us more than just the null space. A coefficient $a_{j}$ of the model produces a contribution 
$s_{j}a_{j}\mathbf{u}_{j}$ to the data, and so if $s_{j}$ is small this contribution is easily lost within the measurement 
errors. Directions in the model space associated with small singular values are, therefore, effectively unresolved even 
though they do not lie within the null space, and any attempt to recover them will amplify the noise in the data. 
Fig.~\ref{fig:4} shows the singular values of the toy problem along with two of the singular vectors: the first is 
concentrated where many rays cross, and so is well constrained, while the other lies almost entirely within a few poorly 
sampled cells.

\begin{figure}[ht]
    \centering    
    \includegraphics[width=\textwidth]{figures/L7F4.pdf}        
    \caption{The non-zero singular values of the toy problem, in decreasing order and relative to the largest (left), 
    together with the singular vectors having the largest singular value (centre) and the 251st (right). There are 280 
    data and 576 blocks, so at most 280 singular values are non-zero and the null space has dimension of at least 296.}
    \label{fig:4}
\end{figure}

\section*{Regularised least squares solutions}
To solve the inverse problem, we might be tempted to apply a simple least squares method, and hence consider the
 following weighted \textbf{least squares misfit}\footnote{Some of you might note that this is proportional to the $\chi^{2}$-statistic.}
\begin{equation}
    J(\mathbf{m}) = \frac{1}{2} \sum_{i=1}^{n} \frac{[(\mathbf{A}\mathbf{m})_{i} - d_{i}]^{2}}{\sigma_{i}^{2}},
    \label{eq:18}
\end{equation}
where the factor of one half has been included for later convenience. If we introduce the diagonal \textbf{covariance matrix}
 $\mathbf{C}=\mathrm{diag}(\sigma_{1}^{2},\dots,\sigma_{n}^{2})$ then the misfit above can be written more concisely as
\begin{equation}
    J(\mathbf{m}) = \frac{1}{2} (\mathbf{A}\mathbf{m} - \mathbf{d})^{T} \mathbf{C}^{-1} (\mathbf{A}\mathbf{m} - \mathbf{d}).
    \label{eq:19}
\end{equation}
This second expression is also valid for a general covariance matrix that accounts for correlations between the different
 measurement errors. The problem with this approach is that there is not a unique model that minimises this chosen misfit.
  Indeed, suppose there exists a non-zero vector $\mathbf{m}_{0}$ lying in the null space of $\mathbf{A}$. It is then clear
   that $J(\mathbf{m} + a\mathbf{m}_{0}) = J(\mathbf{m})$ for all scalars $a$.

A non-ideal but pragmatic way around this problem is to seek a model that not only fits the data well, but possesses
 a property that we see as desirable. For example, we might think that the true velocity structure must be close to that
  of the spherical reference model, and hence prefer models that are small in some sense. A common way to measure the
   ``size'' of the velocity perturbation is through the $L^{2}$ norm
\begin{equation}
    \int_{M} \delta\alpha^{2} \dd^{3}\mathbf{x}.
    \label{eq:20}
\end{equation}
To proceed we substitute our finite-dimensional expansion of the velocity perturbation into this expression to obtain
\begin{equation}
    \int_{M} \delta\alpha^{2} \dd^{3}\mathbf{x} = \sum_{j=1}^{m}\sum_{k=1}^{m} \delta\alpha_{j}\delta\alpha_{k} \int_{M}
     \varphi_{j}\varphi_{k} \dd^{3}\mathbf{x},
    \label{eq:21}
\end{equation}
which in matrix-vector notation can be written
\begin{equation}
    \int_{M} \delta\alpha^{2} \dd^{3}\mathbf{x} = \mathbf{m}^{T}\mathbf{B}\mathbf{m},
    \label{eq:22}
\end{equation}
where the matrix $\mathbf{B}$ has components
\begin{equation}
    B_{jk} = \int_{M} \varphi_{j}\varphi_{k} \dd^{3}\mathbf{x}.
    \label{eq:23}
\end{equation}
Given this result we now consider the modified least squares misfit
\begin{equation}
    J(\mathbf{m}) = \frac{1}{2} (\mathbf{A}\mathbf{m} - \mathbf{d})^{T} \mathbf{C}^{-1} (\mathbf{A}\mathbf{m} -
     \mathbf{d}) + \frac{\lambda}{2} \mathbf{m}^{T}\mathbf{B}\mathbf{m},
    \label{eq:24}
\end{equation}
where $\lambda$ is a positive constant. The first term on the right-hand side measures how well the data are fit, while
 the second is proportional to the quantity we believe will be small for the true model. The \textbf{damping parameter}, $\lambda$,
provides a means of trading off between the two terms. By making $\lambda$ smaller we favour models that fit the data more 
closely, while larger values favour models compatible with our preconceptions. The effect of the damping parameter on the 
toy problem is shown in Fig.~\ref{fig:5}, along with the corresponding trade-off between data misfit and model norm. 
How, then, should $\lambda$ be chosen? A natural criterion is to require that the data are fit to within their errors, 
but no better, so that the misfit takes its expected value $\chi^{2}\approx n$. This is known as the \textbf{discrepancy 
principle}, and it was used to fix the damping in Fig.~\ref{fig:2}; it will be met again in the second problem set. 
Fitting the data more closely than this serves only to fit the noise.

\begin{figure}[ht]
    \centering    
    \includegraphics[width=0.9\textwidth]{figures/L7F5.pdf}        
    \caption{Regularised least squares solutions of the toy problem for three values of the damping parameter, where $\lambda_{*}$ is 
    the value given by the discrepancy principle, along with the trade-off curve of model norm against data misfit 
    as $\lambda$ is varied (the three models are marked). Too small a value of $\lambda$ fits the noise in the data, while too large 
    a value suppresses the model.}
    \label{fig:5}
\end{figure}

Let us show that minimisation of eq.~(\ref{eq:24}) leads to a unique model for each $\lambda > 0$. To do this we will differentiate 
$J$ with respect to $\mathbf{m}$, and show that the equation
\begin{equation}
    \nabla J(\mathbf{m}) = \mathbf{0}
    \label{eq:25}
\end{equation}
has a unique solution corresponding to a minimum of $J$. Using index notation we have
\begin{equation}
    J(\mathbf{m}) = \frac{1}{2}(A_{ij}m_{j}-d_{i})C_{ik}^{-1}(A_{kl}m_{l}-d_{k}) + \frac{\lambda}{2}m_{i}B_{ij}m_{j},
    \label{eq:26}
\end{equation}
and so
\begin{align}
    \frac{\partial J}{\partial m_{p}} =& \frac{1}{2}(A_{ip})C_{ik}^{-1}(A_{kl}m_{l}-d_{k}) \nonumber \\
    &+ \frac{1}{2}(A_{ij}m_{j}-d_{i})C_{ik}^{-1}(A_{kp}) \nonumber \\
    &+ \frac{\lambda}{2} \delta_{ip}B_{ij}m_{j} + \frac{\lambda}{2} m_{i}B_{ij}\delta_{jp} \nonumber \\
    =& A_{pi}^{T}C_{ik}^{-1}(A_{kl}m_{l}-d_{k}) + \lambda B_{pj}m_{j}.
    \label{eq:27}
\end{align}
Reverting to matrix-vector notation this result becomes
\begin{equation}
    \nabla J(\mathbf{m}) = (\mathbf{A}^{T}\mathbf{C}^{-1}\mathbf{A} + \lambda \mathbf{B})\mathbf{m} - \mathbf{A}^{T}\mathbf{C}^{-1}\mathbf{d}.
    \label{eq:28}
\end{equation}
Setting this derivative equal to zero we arrive at the linear equations
\begin{equation}
    (\mathbf{A}^{T}\mathbf{C}^{-1}\mathbf{A} + \lambda \mathbf{B})\mathbf{m} = \mathbf{A}^{T}\mathbf{C}^{-1}\mathbf{d},
    \label{eq:29}
\end{equation}
which will have a unique solution so long as the $m \times m$ symmetric matrix
\begin{equation}
    \mathbf{A}^{T}\mathbf{C}^{-1}\mathbf{A} + \lambda \mathbf{B}
    \label{eq:30}
\end{equation}
is invertible. That this is the case can be seen by examining the sign of its eigenvalues. To do this we consider the eigenvalue problem
\begin{equation}
    (\mathbf{A}^{T}\mathbf{C}^{-1}\mathbf{A} + \lambda \mathbf{B})\mathbf{m} = \mu \mathbf{m},
    \label{eq:31}
\end{equation}
with $\mathbf{m}$ the eigenvector and $\mu$ the corresponding eigenvalue. Taking the inner product of this equation with $\mathbf{m}$ we obtain
\begin{equation}
    (\mathbf{A}\mathbf{m})^{T}\mathbf{C}^{-1}(\mathbf{A}\mathbf{m}) + \lambda \mathbf{m}^{T}\mathbf{B}\mathbf{m} = \mu,
    \label{eq:32}
\end{equation}
where we have assumed that the eigenvector is normalised. The first term on the left-hand side is clearly non-negative, 
while the second is positive by definition. Thus all eigenvalues are positive, and the matrix is invertible. In this 
manner we have shown that the unique solution of $\nabla J(\mathbf{m})=\mathbf{0}$ is given by the simple formula
\begin{equation}
    \mathbf{m} = (\mathbf{A}^{T}\mathbf{C}^{-1}\mathbf{A} + \lambda \mathbf{B})^{-1} \mathbf{A}^{T}\mathbf{C}^{-1}\mathbf{d}.
    \label{eq:33}
\end{equation}
Finally, we note that this model must be a minimum of $J$. This is because it is a quadratic function that tends to 
infinity as $\|\mathbf{m}\|$ grows in magnitude, and hence its only stationary point must be a minimum.

What we have obtained in eq.~(\ref{eq:33}) is known as a \textbf{regularised least squares solution} of the inverse problem. Such a 
solution is obtained by choosing a regularisation matrix $\mathbf{B}$ corresponding to a quantity we wish to make 
small, along with the damping parameter $\lambda$ which determines how data misfit and model size are balanced. 
The action of the damping can be seen clearly using the singular value decomposition of eq.~(\ref{eq:17}). Taking 
$\mathbf{B}=\mathbf{I}$ and $\mathbf{C}=\sigma^{2}\mathbf{I}$ for simplicity, eq.~(\ref{eq:33}) becomes
\begin{equation}
    \mathbf{m}=\sum_{s_{j}>0}\frac{s_{j}}{s_{j}^{2}+\lambda\sigma^{2}}(\mathbf{u}_{j}\cdot\mathbf{d})\,\mathbf{v}_{j},
    \label{eq:34}
\end{equation}
so that the components of the model along well-resolved directions, for which $s_{j}^{2}\gg\lambda\sigma^{2}$, are 
recovered as $(\mathbf{u}_{j}\cdot\mathbf{d})/s_{j}$, exactly as they would be by inverting eq.~(\ref{eq:17}) directly, 
while those along directions with $s_{j}^{2}\ll\lambda\sigma^{2}$ are suppressed, and the null space is never entered. 
Damping thus acts by discarding those parts of the model to which the data are least sensitive.

There is a range of different schemes for regularisation including, for example, those that seek as smooth a model as possible. 
Indeed, a more general least squares function to minimise is
\begin{equation}
    J(\mathbf{m}) = \frac{1}{2}(\mathbf{A}\mathbf{m}-\mathbf{d})^{T}\mathbf{C}^{-1}(\mathbf{A}\mathbf{m}-\mathbf{d})
     + \frac{\lambda}{2}(\mathbf{m}-\mathbf{m}_{0})^{T}\mathbf{B}(\mathbf{m}-\mathbf{m}_{0}),
    \label{eq:35}
\end{equation}
where we can see the regularisation term as asking that the model be close to some preferred value $\mathbf{m}_{0}$, with 
``close'' being measured by the matrix $\mathbf{B}$. Following the same method as above, the optimal solution in this case 
can be shown to be
\begin{equation}
    \mathbf{m} = (\mathbf{A}^{T}\mathbf{C}^{-1}\mathbf{A} + \lambda \mathbf{B})^{-1} (\mathbf{A}^{T}\mathbf{C}^{-1}\mathbf{d} 
    + \lambda \mathbf{B}\mathbf{m}_{0}).
    \label{eq:36}
\end{equation}
Though such approaches are all rather ad hoc, they are widely used because of their simplicity and applicability 
to problems in high-dimensional spaces. 

Whatever the regularisation used, it is important to assess what a given data set is actually able to resolve. A simple 
and instructive test is to invert synthetic data calculated from an input model that is zero except within a single small 
region, a so-called \textbf{spike test}, using the same ray paths and regularisation as for the real data. If the data 
constrained the model perfectly the spike would be recovered exactly; in practice its amplitude is reduced and it is 
smeared out along the ray paths that cross it, to an extent that depends strongly on the ray coverage, as Fig.~\ref{fig:6} 
shows for the toy problem. Because the inversion is linear, the recovered model for a spike is a column of the matrix that 
maps input models to recovered models, and this \textbf{resolution matrix} will be considered further in the second problem set.

\begin{figure}[ht]
    \centering    
    \includegraphics[width=0.85\textwidth]{figures/L7F6.pdf}        
    \caption{Spike tests for the toy problem. Noise-free synthetic data are calculated for an input model that is non-zero only 
    within a single block (left), and are inverted with the regularisation used in Fig.~\ref{fig:2} (right). For a spike in a 
    well-covered block (upper row, crossed by 39 rays) about two thirds of the amplitude is recovered in the right place, with some 
    smearing along the rays that cross it. For a spike in a poorly covered block (lower row, crossed by four rays) only about 
    a fifth of the amplitude is recovered, and the anomaly is spread along the few rays that sample it: such a block is 
    barely constrained by the data, and the damping suppresses it.}
    \label{fig:6}
\end{figure}

\section*{A Bayesian perspective}
A seemingly very different way of thinking about inverse problems is through \textbf{Bayes' theorem}. This states that 
 the \textbf{probability density function} (PDF) of the model, $\mathbf{m}$, given that
the data, $\mathbf{d}$, have been observed can be written
\begin{equation}
    p(\mathbf{m}|\mathbf{d}) = \frac{ p(\mathbf{d}|\mathbf{m}) p(\mathbf{m})}{\int p(\mathbf{d}|\mathbf{m}) p(\mathbf{m}) \dd \mathbf{m}}, 
    \label{eq:37}
\end{equation}
where $p(\mathbf{d}|\mathbf{m})$, which is known as the \textbf{likelihood}, is the PDF for observing $\mathbf{d}$ from the model $\mathbf{m}$, 
and $p(\mathbf{m})$ is the PDF of $\mathbf{m}$ prior to observing any data. In words, this theorem tells us how to update uncertain 
knowledge of the model in light of our observations. This result is fully general, and can be applied to non-linear inverse problems 
with complicated priors. To do so in practice, however, requires random sampling techniques such as the \textbf{Markov chain Monte Carlo} (MCMC) 
algorithm, and this is generally possible only for relatively low-dimensional problems for which forward calculations have a 
low computational cost.

As a special case, suppose we consider a linear inverse problem with Gaussian data errors, and assume a Gaussian 
prior on the model with mean $\mathbf{m}_{0}$ and covariance $\mathbf{C}_{m}$. The prior PDF then takes the form
\begin{equation}
    p(\mathbf{m}) = \frac{1}{\sqrt{(2\pi)^{m}\det \mathbf{C}_{m}}} \exp\left[-\frac{1}{2}(\mathbf{m}-\mathbf{m}_{0})^{T}
    \mathbf{C}_{m}^{-1}(\mathbf{m}-\mathbf{m}_{0})\right],
    \label{eq:38}
\end{equation}
where we recall that $m$ is the number of basis vectors for the model. Similarly with Gaussian errors, the likelihood can be written
\begin{equation}
    p(\mathbf{d}|\mathbf{m}) = \frac{1}{\sqrt{(2\pi)^{n}\det \mathbf{C}}} \exp\left[-\frac{1}{2}(\mathbf{A}\mathbf{m}
    -\mathbf{d})^{T}\mathbf{C}^{-1}(\mathbf{A}\mathbf{m}-\mathbf{d})\right].
    \label{eq:39}
\end{equation}
From Bayes' theorem we then find
\begin{equation}
    p(\mathbf{m}|\mathbf{d}) \propto \exp\left[-\frac{1}{2}(\mathbf{m}-\mathbf{m}_{0})^{T}\mathbf{C}_{m}^{-1}(\mathbf{m}
    -\mathbf{m}_{0}) - \frac{1}{2}(\mathbf{A}\mathbf{m}-\mathbf{d})^{T}\mathbf{C}^{-1}(\mathbf{A}\mathbf{m}-\mathbf{d})\right],
    \label{eq:40}
\end{equation}
and a little effort shows that the exponentiated term can be factored into the form
\begin{equation}
    p(\mathbf{m}|\mathbf{d}) \propto \exp\left[-\frac{1}{2}(\mathbf{m}-\mathbf{m}_{p})^{T}\mathbf{C}_{p}^{-1}
    (\mathbf{m}-\mathbf{m}_{p})\right],
    \label{eq:41}
\end{equation}
where we have defined the \textbf{posterior mean}
\begin{equation}
    \mathbf{m}_{p} = (\mathbf{A}^{T}\mathbf{C}^{-1}\mathbf{A} + \mathbf{C}_{m}^{-1})^{-1} (\mathbf{A}^{T}
    \mathbf{C}^{-1}\mathbf{d} + \mathbf{C}_{m}^{-1}\mathbf{m}_{0}),
    \label{eq:42}
\end{equation}
and the \textbf{posterior model covariance}
\begin{equation}
    \mathbf{C}_{p} = (\mathbf{A}^{T}\mathbf{C}^{-1}\mathbf{A} + \mathbf{C}_{m}^{-1})^{-1}.
    \label{eq:43}
\end{equation}
We conclude that in this special case the posterior distribution is Gaussian, and can be determined in closed form. 
Moreover, the posterior mean is precisely that obtained via the regularised least squares method subject to the 
choice $\lambda \mathbf{B} = \mathbf{C}_{m}^{-1}$. The posterior covariance provides, in addition, a measure of the 
uncertainty in the recovered model. Fig.~\ref{fig:7} shows the posterior mean and standard deviation for the toy problem 
using a Gaussian prior whose covariance between two cells decays with their separation over a correlation length of 
about a sixth of the box, this encoding the expectation that the slowness perturbation is smooth. The posterior mean 
recovers the input model closely, including within regions crossed by few rays, where the prior correlation carries 
information from neighbouring cells, and the posterior standard deviation shows where this recovery can be trusted.

\begin{figure}[ht]
    \centering    
    \includegraphics[width=\textwidth]{figures/L7F7.pdf}        
    \caption{Bayesian solution of the toy problem using a Gaussian prior with a correlation length of $0.15$ times the box side and 
    a standard deviation of $0.003$: the input model (left), the posterior mean (centre), and the posterior standard 
    deviation (right), which approaches that of the prior only in the corners far from any ray.}
    \label{fig:7}
\end{figure} In this manner we see that solving a regularised least squares 
problem returns the mean model corresponding to a certain prior distribution. It is sometimes argued that this 
result justifies the use of regularised least squares, with the prior distribution 
determining the appropriate choice of regularisation. It is better to read the result the other way round. 
Any choice of regularisation amounts to the choice of a prior, but one that is made implicitly, and then adjusted through 
a damping parameter in a rather ad hoc manner. Within the Bayesian approach the assumptions made about the model are stated 
explicitly through the prior, their consequences follow from the rules of probability, and the result is a distribution 
rather than a single model, so that the uncertainty in what has been recovered is quantified. That the prior remains a matter 
of judgement is not a weakness of the method: some such choice is unavoidable in any underdetermined problem, and the Bayesian 
formulation simply makes it honest and visible. 

\subsection*{Inversion in function spaces}
Throughout this lecture the model was reduced at the outset to a finite set of coefficients through the expansion of 
eq.~(\ref{eq:8}), and the choice of basis made there is arbitrary. An alternative, developed over roughly the past fifteen years, is to 
pose and solve the inverse problem for the function $\delta\alpha$ itself. The model is taken to lie within a suitable space of 
functions, one in which smoothness can be measured, the prior becomes a Gaussian probability measure on this space, and the 
data remain finite in number. Bayes' theorem can be formulated in this setting, and for a linear forward problem with Gaussian 
errors the posterior is again a Gaussian measure, with mean and covariance given by expressions that mirror eqs.~(\ref{eq:39}) 
and (\ref{eq:40}). Only once the problem has been solved in this form is it discretised for the purposes of computation, and 
the essential point is that as the discretisation is refined the discrete answers converge to the function-space solution. 
The result is a method whose answers do not depend on an arbitrary parameterisation, and in which the prior states directly 
what is being assumed about the function. You are not expected to understand these ideas in any detail; the essential point 
to take away is that the problem need not be discretised arbitrarily at the outset, and that when a discretisation is introduced 
it approximates a well-defined answer. Fig.~\ref{fig:8} shows a tomographic problem of the kind considered in this lecture 
solved in this way, using the Python package \texttt{pygeoinf}: the true model is a random draw from the prior, and the posterior 
expectation and standard deviation are shown.\footnote{The theory is set out in the review by Stuart (2010).}

\begin{figure}[ht]
    \centering    
    \includegraphics[width=\textwidth]{figures/L7F8.pdf}        
    \caption{A path-average tomography problem posed and solved in a function space: the true model, which is a random draw 
    from the prior (left), the posterior expectation (centre), and the posterior standard deviation (right). Stars mark sources 
    and circles receivers. The calculations have, of course, still been discretised, as they must be on a computer, but the 
    discretisation is of a well-defined problem and has been taken far enough for the results to have converged in practice. 
    The results are therefore independent of the particular discretisation used, in contrast to those of the pixel basis in 
    the earlier figures. Calculated with \texttt{pygeoinf}.}
    \label{fig:8}
\end{figure}


\section*{What you need to know and be able to do}
\begin{enumerate}
    \item[(i)] Why the construction of tomographic models is of interest: it gives the main quantitative 
    constraint on mantle convection and hence on the dynamics of the Earth's deep interior. Here you should 
    be familiar with the idea that seismic slow regions are likely to be hotter than average, while fast 
    regions will be cooler.
    \item[(ii)] How Fermat's principle is used to obtain a simple linearised relation between delay times 
    and velocity perturbations. From this, you should be able to outline how a tomographic problem is reduced 
    to the standard form $\mathbf{d} = \mathbf{A}\mathbf{m} + \mathbf{e}$.
    \item[(iii)] What the null space of a linear(ised) inverse problem is and why it matters.
    \item[(iv)] How regularised least squares is used to obtain solutions of linearised tomographic problems. 
    Here you should be able to derive the least squares equations from first principles.
    \item[(v)] How Bayes' theorem can be applied in the solution of inverse problems, and that the 
    results obtained coincide with regularised least squares for linear(ised) problems with Gaussian errors and priors.
\end{enumerate}
